% Section .............................................................................................................
\section{Bohr's Provisional Resolution and the Persistence of the Problem}
\label{section:bohr}

% Subsection ..........................................................................................................
\subsection{The 1913 Postulates as Ontological Stipulation}

Bohr's response in 1913 was philosophically unusual in its candour \cite{bohr_1913_constitution}. He did not attempt to repair the classical ontology. He declared classical electrodynamics inapplicable to intra-atomic orbits and introduced quantised stationary states as irreducible postulates: the electron in a stationary state does not radiate, not because some mechanism suppresses radiation, but because the very concept of radiating simply does not apply to that state.

\medskip

This is an ontological statement, not a dynamical explanation. Bohr acknowledged as much, and indeed he described his postulates as provisional, awaiting a deeper theory. The immediate success was spectacular. The energy levels
\begin{equation}
  E_n = -\frac{m_e e^4}{2\hbar^2}\cdot\frac{1}{n^2} = -\frac{13.6\,\mathrm{eV}}{n^2}, \qquad n = 1, 2, 3, \ldots
  \label{eq:bohr-levels}
\end{equation}

reproduce the Balmer series \cite{balmer_1885_spectrallinien} and the full Rydberg formula \cite{rydberg_1888_linespectra}, with the Rydberg constant now derived rather than fitted:
\begin{equation}
  R_H = \frac{m_e e^4}{4\pi\hbar^3 c} \approx 1.097\times10^7\,\mathrm{m}^{-1}.
\end{equation}

The precision was stunning. Yet, the ontological justification was absent.

% Subsection ..........................................................................................................
\subsection{Quantum Mechanics: Relocation, Not Dissolution}

Full quantum mechanics relocates rather than dissolves the problem. In the Schr\"{o}dinger picture, the ground-state electron occupies the state $\psi_{1s}(r) \propto e^{-r/a_0}$, spherically symmetric. There is no definite orbit, no classical centripetal acceleration, and therefore no Larmor radiation. Stability is explained: the ground state is the lowest eigenvalue of the Hamiltonian $\hat{H} = \hat{p}^2/2m_e - e^2/4\pi\epsilon_0 r$, and there is no lower state to decay into.

\medskip

Quantum electrodynamics (QED) goes further, correctly computing excited-state lifetimes, the Lamb shift ($\sim$1057\,MHz for the $2s$--$2p_{1/2}$ splitting, measured by Lamb and Retherford in 1947 \cite{lamb_1947_finestructure}), and spontaneous emission rates at parts-per-billion precision.

% Subsection ..........................................................................................................
\subsection{The Tail That Remains Open}

Yet the problem has a residue. The self-energy of the electron in its own electromagnetic field has no fully satisfactory formulation. The Abraham--Lorentz equation \cite{abraham_1903_dynamik, lorentz_1904_electromagnetic, dirac_1938_radiating},
\begin{equation}
  m_e \dot{v} = F_{\mathrm{ext}} + \frac{e^2}{6\pi\epsilon_0 c^3}\ddot{v},
  \label{eq:AL}
\end{equation}
admits \emph{runaway solutions} (acceleration growing exponentially without external force) and \emph{pre-acceleration solutions} (particle responds before the force arrives, violating causality). In QED these pathologies are controlled by renormalisation, which works numerically with extraordinary precision but which Dirac regarded as a fundamental logical inconsistency \cite{dirac_1938_radiating} and Feynman called sweeping the problem under the rug with an elegant technique \cite{feynman_1985_qed}.

The ontological question Bohr implicitly raised in 1913---what \emph{kind of thing} is an electron, considered as both a quantum object and a source of its own field?---remains incompletely answered. Each theoretical layer has provided a more adequate partial ontology. None has provided a complete one.

\medskip

Consider a paragraph on effective field theory (Weinberg, Wilson): EFT explicitly embraces ontological incompleteness at each scale. That may be a more honest framing than ``the problem is unsolved.''
